% =========================
% 3_motivation.tex
% =========================
\section{Motivation}
\label{sec:motivation}

\subsection{Application Requirements}
\label{sec:why}

Rehabilitation, occupational monitoring, and small field studies motivate an open activity set,
several sensor configurations, and useful behavior from few demonstrations. The application, not
the pretraining corpus, determines the relevant labels. Inertial traces are also difficult to label
in isolation: annotation commonly requires synchronized video, a scripted protocol, or self-report.
A memory-based system can accept a small number of demonstrations without fitting a new model and
can keep those recordings local to the deployment.

\subsection{Acquisition Configuration Determines Observability}
\label{sec:no_grammar}
\label{sec:heterogeneity_motiv}

The recorded signal depends jointly on mounting orientation, body placement, sampling rate,
device-side filtering, gravity convention, and sensor availability. Some differences are explicit
in metadata; others are properties of the conversion. For example, one source may expose total
acceleration while another exposes gravity-removed acceleration. Upsampling changes an array's
length but cannot restore bandwidth absent at acquisition. Likewise, full rotation invariance can
discard posture information carried by gravity.

These facts motivate a distinction between \emph{compatibility}, \emph{admissibility}, and
\emph{similarity}. Compatibility asks whether two rows have comparable modalities and gravity
states. Admissibility asks whether both sensor configurations can observe the candidate concept. A
pocket accelerometer can be relevant to gait but may be uninformative for a fine arm gesture.
Similarity ranks the remaining evidence. Applying these operations in the opposite order permits a
geometrically close but physically irrelevant row to occupy a limited retrieval slot.

HALO therefore represents three kinds of acquisition information. Frequency bands are defined in
hertz and carry masks for acquisition bandwidth and patch resolution. One frozen text descriptor
states device, modality, placement, gravity convention, and partner-sensor presence. A separate
activity-invariant vector measures acquisition behavior such as gravity magnitude, noise floor,
quantization, clipping, rate fidelity, and rest bias. Missing measurements have explicit support
bits. This split permits the paper to test whether the representation learned motion or merely
recognized the source dataset.

\subsection{Activity-Name Similarity Is Not a Motion Metric}
\label{sec:geometry}

An open-label model cannot rely on one fixed output head. A common alternative embeds candidate
names and aligns sensor features with the same sentence space~\cite{clip,lanhar,goat,unimts}. That
space reflects linguistic relation, not only observable motion. \emph{Drinking coffee} and
\emph{drinking water} differ in object semantics while a wrist sensor may observe nearly the same
gesture. Conversely, \emph{sitting} and \emph{sitting down} have similar words but denote a posture
and a transition. Figure~\ref{fig:geometry} presents these as conceptual counterexamples; it does
not assign quantitative motion distances by author judgment.

% =====================================================================
% fig_geometry.tex  --  L2 evidence: sensor similarity != linguistic similarity
%
% DATA PROVENANCE: cosine similarities between all-MiniLM-L6-v2 sentence
% embeddings of the label strings, measured 2026-07-26 by
%   sentence_transformers.SentenceTransformer('all-MiniLM-L6-v2')
% with normalize_embeddings=True.  Raw values in the arrays below; the
% generating script and its JSON output are archived with the paper source.
% The MOTION grouping (same vs different physical motion) is our own
% judgement and is stated as such in the caption -- it is not measured.
% =====================================================================
\begin{figure}[t]
\centering
\setlength{\abovecaptionskip}{3pt}
\setlength{\belowcaptionskip}{-4pt}
\resizebox{\columnwidth}{!}{%
\begin{tikzpicture}
\begin{axis}[
  width=\columnwidth, height=4.35cm,
  xmin=0.40, xmax=1.02,
  ymin=-1.15, ymax=1.85,
  xtick={0.4,0.5,0.6,0.7,0.8,0.9,1.0},
  xticklabel style={font=\scriptsize},
  xlabel={cosine similarity of the two label strings (frozen sentence encoder)},
  xlabel style={font=\scriptsize,yshift=2pt},
  ytick={0,1},
  yticklabels={},
  axis y line=none,
  axis x line*=bottom,
  clip=false,
  scale only axis,
]
% ---- shaded overlap region: where the two groups interleave ----
\addplot[draw=none,fill=halogray!12,forget plot]
  coordinates {(0.5239,-1.15) (0.7965,-1.15) (0.7965,1.85) (0.5239,1.85)}
  \closedcycle;
\node[font=\scriptsize\itshape,halogray,anchor=south,align=center]
  at (axis cs:0.660,1.50) {overlap};

% ---- lane guides ----
\addplot[draw=halogray!35,thin,forget plot] coordinates {(0.40,1) (1.02,1)};
\addplot[draw=halogray!35,thin,forget plot] coordinates {(0.40,0) (1.02,0)};

% ---- TOP LANE: DIFFERENT physical motion, near-identical text ----
\addplot[only marks,mark=*,mark size=1.9pt,color=haloorange,forget plot]
  coordinates {(0.5239,1) (0.6396,1) (0.7668,1) (0.7690,1) (0.8678,1) (0.8990,1) (0.9342,1)};
% group mean
\addplot[only marks,mark=|,mark size=5pt,line width=1.1pt,color=haloorange!70!black,forget plot]
  coordinates {(0.7715,1)};

% ---- BOTTOM LANE: SAME physical motion, distant text ----
\addplot[only marks,mark=*,mark size=1.9pt,color=haloblue,forget plot]
  coordinates {(0.4500,0) (0.4595,0) (0.4917,0) (0.5599,0) (0.6645,0) (0.7775,0) (0.7965,0)};
\addplot[only marks,mark=|,mark size=5pt,line width=1.1pt,color=haloblue!70!black,forget plot]
  coordinates {(0.5999,0)};

% ---- lane titles ----
\node[font=\scriptsize\bfseries,color=haloorange!75!black,anchor=west]
  at (axis cs:0.405,1.42) {different motion, near-identical text};
\node[font=\scriptsize\bfseries,color=haloblue!75!black,anchor=west]
  at (axis cs:0.405,-0.72) {same motion, distant text};

% ---- callouts on the two headline pairs ----
\node[font=\scriptsize,anchor=south east,inner sep=1pt]
  (sitdown) at (axis cs:1.015,1.06) {\emph{sitting} / \emph{sitting down}};
\node[font=\scriptsize,anchor=north west,inner sep=1pt]
  (coffee) at (axis cs:0.500,-0.20) {\emph{drinking coffee} / \emph{drinking water}};
\draw[haloblue,thin] (axis cs:0.5599,-0.10) -- (axis cs:0.552,-0.30);

% ---- mean labels ----
\node[font=\scriptsize,color=haloorange!70!black,anchor=south] at (axis cs:0.7715,1.10) {$\bar{s}=.77$};
\node[font=\scriptsize,color=haloblue!70!black,anchor=south]   at (axis cs:0.5999,0.10) {$\bar{s}=.60$};
\end{axis}
\end{tikzpicture}}
\caption{\textbf{Sensor similarity and linguistic similarity are inverted, not merely
noisy.} Each dot is one pair of activity names from our label vocabulary, placed at the
cosine similarity of their frozen sentence embeddings (all-MiniLM-L6-v2, the encoder
language-aligned HAR models use). Pairs in the \emph{top} lane denote physically distinct
activities; pairs in the \emph{bottom} lane denote the same physical motion. If text
similarity tracked sensor similarity the lanes would separate, with the bottom lane on the
right. They separate the wrong way: the mean of the physically \emph{distinct} pairs
($0.77$) exceeds the mean of the physically \emph{identical} pairs ($0.60$), and $38$ of the
$49$ cross-lane comparisons are inverted. The motion grouping is our judgement; the
similarities are measured. A single projection from sensor space into this text geometry is
therefore mis-specified---which is why HALO retrieves in sensor space and votes in label
space (Section~\ref{sec:evidence_design}).}
\label{fig:geometry}
\end{figure}


HALO uses language in two restricted roles: to describe acquisition configuration and to name
candidates or corpus labels. It does not train the encoder to reproduce sentence distances. Stored
recordings are ranked in sensor-feature space, while configuration-conditioned admissibility decides
whether a neighbor may vote for a candidate. An enrolled demonstration is bound to the candidate by
identity, so arbitrary names remain usable when labeled examples are present.

\textbf{Dependency among components.}
\label{sec:combination}%
Frequency coordinates and masks make rows from different sampling regimes interpretable. Sensor
folding and descriptors expose configuration to the encoder. The evidence engine then uses that
configuration to reject or downweight rows whose placement cannot observe a candidate. The thesis is
therefore not that each component is individually novel, but that acquisition metadata calibrates
evidence rather than being normalized away.
